\scanpage{21}

From now on, in (5.1) we are going to confine ourselves to the case of
anabelian pairs $(X,S)$.

\deleted{To such a}

5.4. To a pair $(X,S)$ (not necessarily \unsure{anabelian}) --- more generally to
$\forall$ scheme $X$ locally of finite type over \unsure{$S$|$\xi$}, we associate an
object ``of a Galois nature'', namely the fundamental groupoid
\inserted{(profinite)} $\Pi_1(X)$ of $X$ (formed if we like of the universal
coverings of $X$), \uline{equipped} with the canonical \deleted{\unsure{functor}}
\inserted{\unsure{homomorphism}}
\diagram{
\begin{tikzcd}
\Pi_1(X) \arrow[r] & \Pi_1(K) \arrow[d, "\simeq"'] \\
& {[\operatorname{Tors}(\Gamma)]}
\end{tikzcd}
}

A morphism of $K$-schemes
\[
X' \xrightarrow{f} X
\]
defines a functor
\[
\Pi_1(X') \xrightarrow{\Pi_1(f)} \Pi_1(X)
\]
and a commutation isomorphism $\alpha$

\margin{20}

\scanpage{22}
\mspage{9}

\diagram{
\begin{tikzcd}[column sep=small]
\Pi_1(X') \arrow[rr, "\Pi_1(f)"]
  \arrow[dr, ""{name=L, below}] &&
\Pi_1(X) \arrow[dl, ""{name=R, below}] \\
& \Pi_1(K)
\arrow[from=L, to=R, bend left=25, "\alpha", "\simeq"'] &
\end{tikzcd}
}

We thus find a functor, from the category of schemes
$X$ locally of finite type over $K$, into the ``category of profinite groupoids over
$\Pi_1(K)$'', defined as in n\textsuperscript{o}\,4.

When we pass to the category of connected $K$-schemes locally of finite type,
\deleted{\unsure{and}} \uline{equipped with a geometric point,
over} $\overline K/K$ (i.e. with an $x\in X$,
\deleted{and} a separable closure $\overline{k(x)}$ of $k(x)$ and a
$K$-morphism $\overline K\hookrightarrow\overline{k(x)}$), this corresponds
to a functor from $K$-schemes locally of finite type and connected, pointed over
$\overline K/K$ (in the preceding sense) to the category of
\deleted{\unsure{extensions}} profinite groups \deleted{\illegible{1}} $\pi$
equipped with a homomorphism (of profinite groups)
\[
\pi\longrightarrow\Gamma
\]

\margin{it is better to say: equipped with a universal covering ....}
\margin{21}

\scanpage{23}

(whose \deleted{\unsure{a}}image will moreover necessarily be open hence
of finite index, for objects coming from $X$ as above).

\uline{Conjecture} 5.5. The restriction of the preceding functor
\inserted{$X\longmapsto(\Pi_1(X)\text{ over }\Pi_1(K))$}
to smooth projective schemes of dimension $\leq 1$ and anabelian
(i.e. such that $(X,S)$ is anabelian, where $S$ is the union of the
components of dimension $0$) is fully faithful.

\margin{This is a bit false \quad cf. \S\,9}

It comes to the same thing to say this:

\deleted{The functor}

\uline{Conjecture 5.5 bis} The functor which, to any $X$
\deleted{connected} as in 5.5. and moreover \uline{connected},
(of dimension $0$ or $1$), equipped with a geometric point $\xi$ over
$\overline K$, associates \deleted{\unsure{the profinite extension}} profinite
group \unsure{$\pi_1(X,\xi)$} over
\unsure{$\Gamma=\pi_1(K,\xi)$}), is a

\margin{\unsure{ditto}}
\margin{22}

\scanpage{24}
\mspage{10}

fully faithful functor.

We must all the same make explicit the
morphisms $(X,\xi)\longrightarrow(X',\xi')$ in the source category:
morphisms of $K$-schemes $X\xrightarrow{f}X'$, equipped with a morphism
\inserted{of $\Pi_1(X')$} (\unsure{or|in} class of paths)
\deleted{\unsure{of}} $f(\xi)\simeq\xi'$.

These conjectures reduce to Galois theory, for $X$ of
dimension $0$. For $X$ of dimension $1$, they concern only $X$ such
that the connected components of $X_{\overline K}$ are of genus $\geq 2$
(or, what comes to the same thing, introducing
the finite extension $K'=H^0(X,\underline{O}_X)$ of $K$, so that $X$ is
geometrically connected over $K'$, \unsure{such that} $X$ as
an algebraic \deleted{\unsure{$K'$-scheme}} \inserted{curve}
\uline{over $K'$} is of genus $\geq 2$. We see easily
(\unsure{taking} $X'=\operatorname{Spec}(K)$,
$X$ \deleted{\illegible{1} \unsure{geometric}} $=$
\unsure{$\mathbf P^1_K$|$\mathbb P^1_K$} or an elliptic curve over $K$)
that they become false otherwise ---

\margin{23}

\scanpage{25}

this is why it was necessary to introduce $S$,
\inserted{plus the anabelian hypothesis on $(X,S)$,}
in order to associate to $(X,S)$ a richer structure than $\Pi_1(X)$ over
$\Pi_1(K)$. We shall find conjectures \inserted{\unsure{(for example)}}
for $X$ \deleted{geometric} a geometrically connected curve over $K$ of genus
\deleted{\unsure{$0$}}$1$ (resp. \deleted{\unsure{$1$}}$0$),
\uline{provided} that $S$
\deleted{\unsure{restriction} \unsure{of} $S$ :}
\inserted{is of degree $\geq 1$ (resp. $\geq 3$)}.

6. \uline{Local analysis of} $(X,S)$ \uline{at an} $s\in S$.

We are interested in the case where $\dim_s(X)=1$, i.e. where $s$ is not an
isolated point of $X$.

Let $O_s$ be the henselization (\inserted{or the} completion, if we insist)
of $\underline{O}_{X,s}$, $K_s$ its field of fractions,
\unsure{$D_s^*$}$=\operatorname{Spec}(K_s)$, --- we identify $s$ with
$\operatorname{Spec} k(s)$ ($k(s)$ is the residue field of
the gauge ring $O_s$). Let us consider
$D_s=\operatorname{Spec}(O_s)$, ---
(``arithmetic disk relative \unsure{to|:} $k(s)$''), thus
$D_s^*$ \inserted{$=D_s\setminus\{s\}=$}
(``punctured disk'') $\longrightarrow D_s$,

\margin{24}

\scanpage{26}
\mspage{11}

we have

(6.1)
\diagram{
\begin{tikzcd}[column sep=large, row sep=large]
\Pi_1(D_s^*)
 \arrow[r, "{\substack{\text{surjective on Hom}\\i_s}}"]
 \arrow[rr, bend left=32, "{\text{\unsure{$p_s$|$\rho_s$}}}"]
 \arrow[dr, dashed, "{\text{\deleted{\unsure{$p_s$|$\rho_s$}}}}", "\sigma_s"']
& \Pi_1(D_s) \arrow[r, "{\substack{\text{faithful}\\q_s}}"]
& \Pi_1(K) \\
& \Pi_1(\text{\deleted{\illegible{1}}}s)
 \arrow[u, "\simeq"]
 \arrow[ur, "j_s"'] &
\end{tikzcd}
}
,

for the choice of a geometric point $\xi_s$ of $D_s^*$
\inserted{over $\overline K/K$} (i.e. of an algebraic closure
$\overline K_s$ of $K_s$ and of a \inserted{$K$-}injection
\unsure{$\overline K\hookrightarrow\overline K_s$|$\overline K\subset\overline K_s$}),
this diagram of groupoids is reflected in a homomorphism of groups

(6.2)
\diagram{
\begin{tikzcd}[column sep=large]
\pi_1(D_s^*\text{\deleted{\unsure{$\xi$}}},\xi_s) \arrow[r, "\text{surjective}"]
 \arrow[d, phantom, "\wr" description]
& \pi_1(k(s),\xi_s) \arrow[r, hook, "\text{injective}"]
 \arrow[d, phantom, "\wr" description]
& \Gamma\text{\deleted{\illegible{1}}} \\
\operatorname{Gal}(\overline K_s/K_s)
& \operatorname{Gal}(\overline{k(s)}/k(s)) &
\end{tikzcd}
}

\margin{\uline{N.B.} the choice of $\overline K_s$ implies a choice of
$\overline k_s$ --- this is the dotted arrow in (6.1)}

whose kernel, as we know from Kummer, is canonically isomorphic to
\unsure{$T_s(\overline{k_s})$|$T_\varepsilon(\overline{k_s})$}
$\simeq T(\overline K_s)$ $[\simeq T(\overline K)]$.
We want to express the datum of this privileged isomorphism as an additional
structure on (6.1) --- \deleted{\unsure{groupoid}}
i.e.; on the groupoid $\Pi_1(D_s^*)$ over $\Pi_1(s)$
(\unsure{i.e.|or} over $\Pi_1(K)$) ---.

We can say it as follows: \deleted{\unsure{In}} If to every
\deleted{\unsure{$\xi$}} $\xi\in\Pi_1(D_s^*)$, we associate the

\margin{25}

\scanpage{27}

kernel of
\[
\operatorname{Aut}(\xi)\longrightarrow\operatorname{Aut}(i(\xi))
\]
(which is also the kernel of the composite
\[
\operatorname{Aut}(\xi)\longrightarrow\operatorname{Aut}(i(\xi))
\longrightarrow
\operatorname{Aut}(\text{\unsure{$p_s$|$\rho_s$}}(\xi)\,
\text{\deleted{\illegible{1}}}=q_s(i_s(\xi)))
\]
we find an \uline{abelian} group, which depends (up to canonical isomorphism)
only \deleted{and forms} on $i(\xi)=\xi'$
[this, \inserted{and the rest of the sentence,}
works whenever we have
\margin{[black right-angle mark, $\urcorner$ or 7]}
a functor
\deleted{\unsure{is surjective}} of \inserted{connected} groupoids,
with abelian kernels \inserted{and surjective on Hom}],
\deleted{\unsure{where}}
and for varying $\xi'$ forms a local system on $\Pi_1(D_s)$,
\deleted{That said, we have}
\inserted{\unsure{can call the} \deleted{\unsure{$\pi_1$}}
\uline{relative} $\pi_1$}
of the groupoid $\Pi_1(D_s^*)$ over \unsure{the the} groupoid
$\Pi_1(D_s)$. That said, we have an isomorphism of local systems of groups
\[
\pi_1\bigl(\Pi_1(D_s^*)\text{ over }\Pi_1(D_s)\bigr)
\simeq q_s^*(\mathcal T_K)
\]
where $\mathcal T_K$ is the Tate local system on $K$.

\margin{\unsure{11|''}}
\margin{26}

\scanpage{28}
\mspage{12}

Let us now set
\[
D_S=\coprod_{s\in S}D_s
\qquad\text{(``\inserted{\unsure{arithmetic}} multidisk at $S$'')}
\]
\[
D_S^*=\coprod_{s\in S}D_s^*
\qquad\text{(``\inserted{\unsure{arithmetic}} multiannulus at $S$'')}
\]

We have a homomorphism of groupoids
\[
\Pi_1(D_S^*)
\xrightarrow{\text{\deleted{\illegible{1}}}\,\sigma_S}
\Pi_1(S)
\quad\bigl(\xrightarrow{j_S}\Pi_1(K)\bigr)
\]
and a canonical isomorphism
\[
\pi_1\bigl(\Pi_1(D_S)/\Pi_1(S)\bigr)
\simeq
j_S^*(\mathcal T(K))
\]

With this set up, we also have a morphism
\[
D_S^*\xrightarrow{\rho_S}\text{\unsure{$X\setminus S$}}
\]
inducing
\[
\Pi_1(D_S^*)\xrightarrow{\rho_S}\Pi_1(X\setminus S).
\]

7. Thus, to $(X,S)$ as in 5.1, we associate:

\deleted{
$1^\circ)$ A (profinite) groupoid
$\Pi_U=\Pi_1(\overbrace{X\setminus S}^{\text{\inserted{$U$}}})$
over $\Pi_e=\Pi_1(K)$ $(\simeq\operatorname{Tors}(\Gamma))$.
}

\margin{\diagram{
\begin{tikzcd}
\Pi_U\arrow[r,"\varphi"]&\Pi_0
\end{tikzcd}
}}
\margin{This groupoid \unsure{$\Pi_0$|$\Pi_1$} \illegible{5} of the connected
components \unsure{which are those} of $X\setminus S$, \unsure{or again}
of $X$}
\margin{27}

\scanpage{29}

\deleted{
$2^\circ)$ A (profinite) groupoid $\Pi_D=\Pi_1(S)$ over $\Pi_0$,
such that $\Pi_D$ has only a finite number of connected components
(\unsure{corresponding} to the points of $S$) and that the functor
$\Pi_D\longrightarrow\Pi_0$ is faithful
}

\margin{\diagram{
\begin{tikzcd}
\Pi_D\arrow[r,"\psi"]&\Pi_0
\end{tikzcd}
}}

\deleted{
$3^\circ)$ A (profinite) groupoid $\Pi_{D^*}=\Pi_1(D_S^*)$,
\deleted{above, over $\Pi_D$, i.e. equipped with a functor}
\deleted{\unsure{$\Pi_D$}}
}

\deleted{
$4^\circ)$ Two functors
\diagram{
\begin{tikzcd}[column sep=small]
& \Pi_{D^*}\arrow[dl,"\rho"']\arrow[dr,"\sigma"] & \\
\text{\unsure{$\Pi_U$}} & \text{//} & \Pi_D
\end{tikzcd}
}
where $\sigma$ is faithful, $\rho$ surjective on Hom and with abelian kernel
}

\deleted{
$5^\circ)$ A commutation isomorphism $\alpha$
\diagram{[Diamond: at the top, $\Pi_{D^*}$; on the left,
\unsure{$\Pi_U$|$\Pi_D$}; on the right,
\unsure{$\Pi_U$|$\Pi_D$}; at the bottom, $\Pi_e$.
Arrow from the top to the left, $\rho$; from the top to the right,
\unsure{$\sigma$|$\rho$}, overwritten; from the left to the bottom,
$\varphi$; from the right to the bottom, $\psi$, with
\deleted{\illegible{1}}. In the center, curved arrow to the right,
bearing $\alpha$ and $\simeq$.]}
}

\margin{28}

\scanpage{30}
\mspage{13}

\unsure{\deleted{$6^\circ)$ A}|\uline{$6^\circ)$ A}}

$1^\circ)$ Three (profinite) groupoids $\Pi_U$, $\Pi_D$, $\Pi_{D^*}$
(in addition to $\Pi_e=\Pi_1(\operatorname{Spec}(K))$
$(\simeq\operatorname{Tors}(\Gamma))$)

$2^\circ)$ \deleted{Two} \inserted{Four} functors
(of profinite groupoids)
\diagram{
\begin{tikzcd}[column sep=small]
& \Pi_{D^*}\arrow[dl,"\rho"']\arrow[dr,"\sigma"] & \\
\Pi_U\arrow[dr,"\varphi"'] && \Pi_D\arrow[dl,"\psi"] \\
& \Pi_e &
\end{tikzcd}
}

$3^\circ)$ A commutation isomorphism
\[
\alpha:\varphi\rho\simeq\psi\sigma
\]
(which is even the identity in the case of the system coming from $(X,S)$,
but \deleted{\illegible{1}} it is better to forget that this is so).

These data satisfy the \deleted{\illegible{1}}
preliminary conditions

a) $\sigma$ induces an isomorphism on the $\pi_0$, and epimorphisms on Hom,
and it has abelian kernel.

b) $\psi$ is faithful

[ c) $\rho$ is faithful ...

d) $\varphi$ is epimorphic modulo finite groups on Aut ... ]

\margin{29}